\chapter{Experimental Validation}
\label{ch:validation}

\section{Experimental Setup}

\subsection{Hardware Configuration}

\begin{table}[htbp]
\centering
\caption{Experimental Hardware}
\begin{tabular}{@{}lll@{}}
\toprule
\textbf{Device} & \textbf{Role} & \textbf{Specifications} \\
\midrule
Smartphone A & Hotspot Host & Android 12, Snapdragon 750G, 6GB RAM \\
Laptop B & Client 1 & Ubuntu 22.04, Intel i5, 8GB RAM \\
Laptop C & Client 2 & Windows 11, Intel i7, 16GB RAM \\
Smartphone D & Client 3 & Android 13, Snapdragon 888, 8GB RAM \\
\bottomrule
\end{tabular}
\end{table}

\subsection{Network Configuration}

\begin{itemize}
    \item \textbf{Cellular Network:} 4G LTE (Carrier: Viettel)
    \item \textbf{Uplink Capacity:} 8 Mbps (measured via speedtest)
    \item \textbf{Downlink Capacity:} 45 Mbps
    \item \textbf{Hotspot Subnet:} 192.168.43.0/24
\end{itemize}

\section{Test Scenarios}

\subsection{Scenario 1: Baseline (No QoS)}

\textbf{Objective:} Establish baseline performance without \qos enforcement.

\textbf{Procedure:}
\begin{enumerate}
    \item Disable \qos scheduler
    \item Run iPerf3 from Client 1 and Client 2 simultaneously
    \item Measure throughput, latency, jitter for 60 seconds
\end{enumerate}

\subsection{Scenario 2: Single HIGH-Priority Device}

\textbf{Objective:} Measure throughput advantage of HIGH-priority device.

\textbf{Procedure:}
\begin{enumerate}
    \item Enable \qos scheduler
    \item Assign Client 1: HIGH, Client 2: LOW
    \item Run iPerf3 from both clients simultaneously
    \item Measure throughput ratio (HIGH/LOW)
\end{enumerate}

\subsection{Scenario 3: Mixed Priorities Under Congestion}

\textbf{Objective:} Validate \qos enforcement under competitive load.

\textbf{Procedure:}
\begin{enumerate}
    \item Enable \qos scheduler
    \item Assign Client 1: HIGH, Client 2: MEDIUM, Client 3: LOW
    \item Run iPerf3 from all clients simultaneously
    \item Measure throughput distribution
\end{enumerate}

\subsection{Scenario 4: Priority Override During Transfer}

\textbf{Objective:} Test dynamic priority changes.

\textbf{Procedure:}
\begin{enumerate}
    \item Start iPerf3 transfer from Client 1 (MEDIUM priority)
    \item After 30 seconds, change Client 1 to HIGH priority
    \item Observe throughput change
\end{enumerate}

\section{Measurement Procedures}

\subsection{Throughput Measurement}

iPerf3 command:
\begin{verbatim}
iperf3 -c <server_ip> -t 60 -i 1 -J > results.json
\end{verbatim}

\subsection{Latency and Jitter Measurement}

Continuous ping:
\begin{verbatim}
ping -i 0.1 -c 600 8.8.8.8 > latency.txt
\end{verbatim}

\subsection{Packet Loss Measurement}

iPerf3 UDP test:
\begin{verbatim}
iperf3 -c <server_ip> -u -b 5M -t 60 > udp_results.txt
\end{verbatim}

\subsection{Performance Profiling}

Android Profiler captures:
\begin{itemize}
    \item CPU usage (percentage)
    \item Memory allocation (MB)
    \item Network I/O (bytes/second)
\end{itemize}

\section{Data Collection}

Each scenario is repeated 10 times to ensure statistical significance. Raw data includes:
\begin{itemize}
    \item iPerf3 JSON output (throughput per second)
    \item Ping logs (RTT per packet)
    \item Android Profiler traces (CPU, memory)
    \item Wireshark packet captures (classification verification)
\end{itemize}
